\documentclass[10pt,a4paper]{amsart}
\usepackage[margin=19mm]{geometry}
\usepackage{amsmath,amssymb,mathtools}
\usepackage[scr=rsfs,cal=euler]{mathalfa}
\usepackage{xfrac}
\usepackage[unicode,hidelinks,bookmarks=false]{hyperref}
\newcommand{\ie}{i.e.\ }
\newcommand{\cf}{cf.\ }
\newcommand{\resp}{respectively\ }
\newcommand{\Subsection}{\subsection}
\providecommand{\nobreakdash}{\nobreak\hbox{-}}
\setlength{\emergencystretch}{2em}
\multlinegap0pt
\usepackage{amssymb}
\newtheorem{lem}{Lemma}
\def\R{\mathbb{R}}
\def\wt{\widetilde}
\title{Faber-Krahn inequalities for first Dirichlet eigenvalues of combinatorial $p$-Laplacian on graphs with boundary}
\author{Wankai He, Chengjie Yu}
\date{Source: arXiv:2603.20814v2 (CC BY 4.0)}
\begin{document}
\maketitle

\noindent\emph{Source and licence.} Faber-Krahn inequalities for first Dirichlet eigenvalues of combinatorial $p$-Laplacian on graphs with boundary. Version arXiv:2603.20814v2 (CC BY 4.0), distributed by arXiv under the Creative Commons Attribution 4.0 International licence, http://creativecommons.org/licenses/by/4.0/, as recorded in the arXiv licence metadata for this exact version and shown on the abstract page https://arxiv.org/abs/2603.20814v2. Full licence notice: \texttt{licenses/arxiv-2603.20814-CC-BY-4.0.txt}.\par\smallskip

\noindent\textit{Editorial scope.} Counted unit: the article's lem lem-comparison-tadpole (source0.tex lines 278--281). The article's complete original proof of it is delivered with it (source0.tex lines 282--340, one \texttt{proof} environment of 59 lines). No mathematics is written, repaired or re-derived: the statement and the proof are the article's own lines, in the article's own notation, and no retained line is substituted or re-flowed.  Retained uncounted as the source-local context the proof needs: lines 202--204; lines 245--247.  Nothing was omitted from any retained span, so the package carries no omission record.  No retained line contains a citation, so the wrapper carries no bibliography and no bibliography entry.  No retained line contains a figure environment, an image-inclusion command or drawing code, so no image is delivered and the package declares no asset dependency.  Editorial formatting only, in addition: an AMS \texttt{amsart} preamble that restates the article's own declarations and macros used by the retained lines, namely the environments \texttt{lem} on separate counters, together with the standard packages those lines need.  The retained environments print in this document as Lemma 1 in order of appearance, whereas the article prints the same items with the numbering of its own full document.  The complete native source of the article as distributed in the arXiv e-print for this exact version is bundled unchanged as \texttt{sources/196880/source0.tex}.
\begin{equation}\label{eq-Rayleigh}
\lambda_{1,p}(G)=\min_{f\in C_B(G)\setminus\{0\}}R_{p,G}[f]
\end{equation}

\begin{lem}\label{lem-max-in-head}
 Let $n>i\geq 3$, $p>1$ and $f$ be a positive first $p$-Dirichlet eigenfunction of $T_{n,i}$. Then, $f$ does not achieve its maximum on the tail of $T_{n,i}$.
\end{lem}

\begin{lem}\label{lem-comparison-tadpole}
For $n\geq 5$ and $p>1$,
$$\lambda_{1,p}(T_{n,4})>\lambda_{1,p}(T_{n,3}).$$
\end{lem}

\begin{proof}
Let $f$ be a positive first $p$-Dirichlet eigenfunction of
$$T_{n,4}: v_n\sim v_{n-1}\sim \cdots\sim v_4\sim v_3\sim v_2\sim v_1\sim v_4.$$
By symmetry and that $\lambda_{1,p}(T_{n,4})$ is of multiplicity one, we know that
\begin{equation}\label{eq-f-v1=v3}
f(v_1)=f(v_3).
\end{equation}
Let $m$ be a maximum vertex of $f$. By Lemma \ref{lem-max-in-head}, $m$ is not in the tail of $T_{n,4}$. So, without loss of generality, $m=v_3$ or $m=v_2$. Let
$$T_{n,3}: u_n\sim u_{n-1}\sim\cdots\sim u_{3}\sim u_2\sim u_1\sim u_3.$$

When $m=v_3$, let $\wt f\in \R^{V(T_{n,3})}$ be such that
\begin{equation*}
\wt f(u_k)=\left\{\begin{array}{ll}f(v_k)&3\leq k\leq n\\
f(m)&1\leq k<3.
\end{array}\right.
\end{equation*}
Then,
\begin{equation*}
\|f\|_{p,T_{n,4}}^p=\sum_{k=1}^nf^p(v_k)\leq \sum_{k=1}^n\wt f^p(u_k)=\|\wt f\|_{p,T_{n,3}}^p
\end{equation*}
and
\begin{equation*}
\begin{split}
\|df\|_{p,T_{n,4}}^p=&\sum_{k=1}^{n-1}|f(v_{k+1})-f(v_k)|^p+|f(v_1)-f(v_4)|^p\\
\geq&\sum_{k=3}^{n-1}|f(v_{k+1})-f(v_k)|^p\\
=&\|d\wt f\|_{p,T_{n,3}}^p.
\end{split}
\end{equation*}
Hence, by \eqref{eq-Rayleigh},
$$\lambda_{1,p}(T_{n,4})=R_{p,T_{n,4}}[f]\geq R_{p,T_{n,3}}[\wt f]>\lambda_{1,p}(T_{n,3}).$$
The last inequality is strict because $\wt f$ is not a first $p$-Dirichlet eigenfunction for $T_{n,3}$ by Lemma \ref{lem-max-in-head}.

When $m=v_2$, let $\wt f\in \R^{V(T_{n,3})}$ be such that
\begin{equation*}
\wt f(u_k)=f(v_k)
\end{equation*}
for $k=1,2,\cdots,n$. By that
\begin{equation*}
\begin{split}
0<&\lambda_{1,p}(T_{n,4})f^{p-1}(v_1)\\
=&\Delta_{p,T_{n,4}} f(v_1)\\
=&|f(v_1)-f(v_4)|^{p-2}(f(v_1)-f(v_4))+|f(v_1)-f(v_2)|^{p-2}(f(v_1)-f(v_2))\\
\leq& |f(v_1)-f(v_4)|^{p-2}(f(v_1)-f(v_4))
\end{split}
\end{equation*}
we have $f(v_1)-f(v_4)>0$. Moreover, by \eqref{eq-f-v1=v3}, 
\begin{equation*}
\begin{split}
\|df\|_{p,T_{n,4}}^p=&\sum_{k=1}^{n-1}|f(v_{k+1})-f(v_k)|^{p}+|f(v_1)-f(v_4)|^p\\
>&\sum_{k=1}^{n-1}|f(v_{k+1})-f(v_k)|^p+|f(v_3)-f(v_1)|^p\\
=&\|d\wt f\|_{p,T_{n,3}}^p,\\
\end{split}
\end{equation*}
and it is clear that
$$\|f\|_{p,T_{n,4}}^p=\|\wt f\|_{p,T_{n,3}}^p.$$
So, by \eqref{eq-Rayleigh},
$$\lambda_{1,p}(T_{n,4})=R_{p,T_{n,4}}[f]>R_{p,T_{n,3}}[\wt f]\geq \lambda_{1,p}(T_{n,3}).$$
This completes the proof of the lemma.
\end{proof}

\end{document}
